\documentclass{article}

\usepackage{maa-monthly}
\usepackage[utf8]{inputenc}
\usepackage[T1]{fontenc}
\usepackage{microtype}

\setlength{\parindent}{0pt}
\setlength{\parskip}{0.75em}

\begin{document}
\thispagestyle{empty}

\begin{center}
{\Large\bfseries Author's Contributions}

\vspace{0.75em}
{\large\itshape The Continuous-Coefficient Jensen Equation:
How Real Mixing Weights Retire Regularity}

\vspace{1em}
{\normalsize Ali Elouafiq}\\
{\small SQLI Digital Lab, Hay Riad, Rabat, Morocco}\\
{\small \texttt{aelouafiq@sqli.com}}
\end{center}

\vspace{1.5em}

\noindent\textbf{Conceptualization.}
A.E.\ conceived the central distinction between the discrete-coefficient
and continuous-coefficient Jensen equations as a separable dictionary
problem, identified the five classical regularity hypotheses that become
vestigial under the continuous-coefficient equation, and formulated the
diagnostic framework (Section~4 of the manuscript) for classifying
saturated-Jensen derivations into four structurally distinct sources.

\medskip
\noindent\textbf{Formal results and proofs.}
A.E.\ formulated and proved all results in the manuscript: the main
theorem (Theorem~1, endpoint substitution), its strict-minimum variant
(Theorem~2), the extension to convex domains of arbitrary real vector
spaces (Theorem~3), the derived regularity corollary (Corollary~1),
the piecewise-saturation corollary (Corollary~2), the irrational-weights
remark (Remark~1), the Hamel-pathology proposition for the
rational-coefficient equation (Proposition~1), and the
resolution-blindness characterization (Proposition~2).

\medskip
\noindent\textbf{Literature review and attribution.}
A.E.\ relied heavily on Acz\'el's corpus, Kuczma's corpus and
bibliography, and Daniel Reem's 2017 bibliography when assembling the
classical regularity-hypothesis literature from Cauchy (1821) through
Kestelman (1947); he also assembled and verified the historical
attribution table (Table~1 and the mirrored Table~2) and identified the
structural parallel with Fischer's 1972 regularity-free entropy
characterization for $n > 2$ outcomes.

\medskip
\noindent\textbf{Formal verification.}
The core lemma of Theorem~1 was initially discovered through a Lean~4
formalization project in which the boundedness hypothesis appeared
unused; A.E.\ designed the formalization harness, interpreted the
formal proof output, and drew the mathematical conclusion presented in
the provenance note.

\medskip
\noindent\textbf{Figures and tables.}
A.E.\ designed and implemented all three manuscript figures in TikZ:
the load diagram (Figure~1, illustrating the dictionary of vestigial
struts), the mechanism diagram (Figure~2, illustrating the endpoint
substitution and the weight axis), and the resolution-axis diagram
(Figure~3, illustrating surrogate calibration and the tent
representation). A.E.\ composed both tables.

\medskip
\noindent\textbf{Writing.}
A.E.\ wrote the manuscript in its entirety, including all sections,
proofs, figure captions, backmatter, and bibliography, and revised it.

\medskip
\noindent\textbf{Submission preparation.}
A.E.\ assembled the submission package, including the full-author and
anonymous manuscript sources, cover letter, and all supplementary
assets, in accordance with the Monthly's double-blind submission policy.

\end{document}